\section{Beam thickness along \texorpdfstring{$z$}{z}}\label{sec:depth}
As a supporting geometric estimate, we fit the angular opacity maps directly for vertical thickness $h=z_{\rm top}-z_{\rm bottom}$. Each beam is a rectangular box with shared bottom $z_0$, width $w$ and thickness $h$. Oblique path lengths depend on thickness as well as width, providing information beyond the image's apparent voxel extent.

The continued-array fit fixes relative beam centres and detector pose, but fits a common shift, bottom, width, thickness and independent nonnegative column opacities. It averages transmitted flux over the admitted aperture before taking its logarithm, and refits a polynomial background. Its conditional values are $h=\ArraydepthH\pm\ArraydepthHSigma\,\mathrm{m}$, $w=\ArraydepthW\pm\ArraydepthWSigma\,\mathrm{m}$ and $z_0=\ArraydepthBottom\pm\ArraydepthBottomSigma\,\mathrm{m}$.

The spreads come from \ArraydepthNReplicas{} count replicas at fixed array and pose. The on-site thickness is $\InputsMeasuredBeamDepth\,\mathrm{m}$; bottom-height comparison with $\InputsMeasuredBeamBottom\,\mathrm{m}$ additionally requires height registration.

This is a conditional estimate: centre, array and background assumptions shift thickness more than its count spread, and a flexible matched-beam fit reaches a thin-depth bound. Noiseless numerical recovery and aperture refinement check the calculation, not real-ceiling accuracy. The weighted residual statistic $\chi^2/\nu=\ArraydepthChisqPerDof$ is a model/weight diagnostic rather than a calibrated goodness-of-fit test.
